% !TEX program = pdflatex
\documentclass[11pt]{article}
\usepackage{support-notes}
\title{Orbit closures on representation spaces detect bounded profiles}
\author{\workingauthor}\date{\coursedate}
\begin{document}\maketitle
\begin{abstract}
In characteristic zero an arbitrary automorphism of the free
localization chart admits a geometric profile test. Compute the generic
rank of a matrix on each component of the orbit closure of a simple
representation. This gives both rows and their order, without periodicity,
including for nonlinear automorphisms. For linear twists in positive
characteristic the corresponding return sets are automatic. The
geometric test fails over $\mathbb F_p(t)$, even in a strict two-term
family. All assertions stated as results are proved.
\end{abstract}

\section{The geometric test}
Use hypotheses (1)--(3) of \othernote{normal_localisations}{Bounded
profiles and finite-dimensional points of universal localisations},
allowing the $k$-algebra automorphism $\sigma$ of
$R=k\langle U_1,\ldots,U_g\rangle$ to have arbitrary order. In
particular $A$ is positive, locally finite and connected, $h$ is
regular normal, the closed fibre forces at least one factor to have
eventually bijective $h$ in every bounded perfect tensor, and all
finite simple points admit perfect lifts on the indicated side.
Call a factor with this eventual bijectivity \emph{tame} in this note.
For a finite-dimensional module choose a basis, giving a point $w$ of
$X_d=(\operatorname{Mat}_d)^g$, $d=\dim_k W$. Substitution by
$\sigma$ gives a polynomial automorphism $\phi$ of $X_d$ with inverse
induced by $\sigma^{-1}$. Use the row convention for right modules and
the column convention for left modules. On the right a basis of row
vectors acts by matrices on the right; the block map on Hom uses
precomposition. With either convention evaluation is polynomial and
equivariant for change of basis.

Let $Z_w$ be the reduced Zariski closure, over an algebraic closure
of $k$, of $\{\phi^j(w):j\in\mathbb Z\}$. For each square matrix $H$
over $R$, write $D_\epsilon(H)$ for the set of isomorphism types of
finite-dimensional simple modules on side $\epsilon=r,l$ on which
$\det H$ is generically nonzero on \emph{every} irreducible component
of $Z_w$. Evaluations use the indicated side. In particular we test a
$k$-simple in its given $k$-dimension; extending the determinant to
the algebraic closure does not require that the simple remain simple
there. For a finite set $F$ of right simple types put
\[
 F^\sharp=\bigcap_{H:\ F\subseteq D_r(H)}D_r(H).
\]
The intersection ranges over square polynomial matrices of all finite
sizes. It contains $F$, and the zero matrix gives $\varnothing^\sharp
=\varnothing$. This definition uses the affine representation spaces
of $R$ and has no quantifier over dg objects.

\begin{thm}\label{geo}
Over a field of characteristic zero all upper-row profiles are labelled,
without duplication, by $F^\sharp$ for finite $F$, ordered by reverse
inclusion. A tame detector supplies the simple factors of its
degree-zero localized cohomology as $F$. All remaining profiles are
labelled by $D_l(H)$, ordered by inclusion. A source $P$ belongs to
this latter profile exactly when $P^\vee$ is tame and every simple
factor of $H^0(P^\vee)_h$ lies in $D_l(H)$. For $g\ge2$ the mixed
order is the same as for a finite twist: no upper-row element lies
below a lower-row element, and a lower-row element lies below an
upper-row element only at the bottom or top endpoints. If $\sigma$
has finite order, $F^\sharp$ is exactly the union of the twist orbits
of $F$ and the finite-order theorem is recovered.
\end{thm}

\begin{lem}\label{return}
For a polynomial automorphism $\phi$ of affine space in characteristic
zero, a polynomial $f$ is nonzero at all but finitely many positive
iterates of $w$ iff it is nonzero generically on each component of
$Z_w$. The same condition governs negative iterates. The exceptional
indices are eventually a finite union of arithmetic progressions.
\end{lem}
\begin{proof}
Both forward and backward tails are dense in $Z_w$. Indeed the closures
of the successive forward tails form a stationary descending chain.
The stable closure is invariant under $\phi$ and its inverse, and contains
the whole orbit. The same proof applies to $\phi^{-1}$ and to any power.
Every component is met infinitely often counting multiplicities of
times: every tail meets its open part outside the other components.
For a periodic point the lemma is immediate. Otherwise the orbit is
injectively parametrized by the integers. The dynamical Mordell--Lang
theorem for \'etale maps of
\href{https://arxiv.org/abs/0808.3266}{Bell--Ghioca--Tucker, Theorem 1.3}
gives a finite union of arithmetic progressions and a finite set of
return times to $V(f)$. Its use over general characteristic-zero
fields is legitimate: the finitely many coefficients of the maps,
$f$ and the point descend to a finitely generated field over $\mathbb Q$,
which embeds in $\mathbb C$. Equalities at all iterates are preserved.
If the returns are infinite, the zero set contains the orbit along
some full progression $a+m\mathbb N$. Its closure contains a component
of $Z_w$: the $m$ translates of the closure of the step-$m$ orbit cover
$Z_w$, so one contains an irreducible component, and translation is
an automorphism permuting the components. Every translate therefore
contains a component. Conversely vanishing identically on a component
forces infinitely many returns by tail density. Repeat for $\phi^{-1}$.
\end{proof}

\begin{proof}[Proof of Theorem~\ref{geo}]
The two local lemmas in the normal-localisation note use no periodicity.
For a tame detector with chart module $W$, present the localized source
by a cone on a finite free matrix. In high degrees its Hom cone has
finite total cohomology exactly when the matrix is bijective at all
but finitely many twists of $W$. This follows degree by degree from
the kernel--cokernel sequence of the cone; a shift of the index leaves
eventuality unchanged. For $W\ne0$ the matrix must be square.
A composition series is preserved by all entries, so its determinant
on $W$ is the product of its determinants on the simple factors.
There are finitely many factors. By Lemma~\ref{return} the upper
test for a square $H$ is exactly $F\subseteq D_r(H)$. Every such
matrix occurs as a test: multiply its entries on one side by a common
power of $h$ to clear denominators, then take its cone. The extra
unit after localization effects a shift or twist, which preserves
the tail test. This also works for noninjective matrices. A rectangular
unequal-size matrix fails on every nonzero $W$ and the case $W=0$
gives the top. Thus containment of upper profiles from $F,F'$ is
equivalent to the implication on all matrix tests, equivalently to
reverse containment of $F^\sharp,(F')^\sharp$. All finite $F$ are
realised by sums of lifts.

For other detectors the closed test forces $P^\vee$ tame. Present the
localized detector by a matrix, tensor with the induced module from
$H^0(P^\vee)_h$, and apply the same determinant argument on the left.
All domains are obtained by adjoining $G=A/(h)$ to a cleared matrix
cone. Duals of simple lifts distinguish these domains. The source
$G$ belongs to all upper and no lower profiles; if finite chart modules
on opposite sides both have positive dimension, the free resolution
of the right one has ranks $gd$ and $d$. For $g\ge2$ its evaluations
have unequal dimensions in every degree, so no nontrivial mixed
inclusion is possible. For finite $\sigma$, the component tests are
exactly evaluations at each twist. The Chinese remainder argument of
the finite-order theorem separates any additional simple orbit from
the finitely many prescribed orbits, giving precisely the stated closure.
\end{proof}

The theorem also describes the lower labels as follows. For each
component pass to its function field and evaluate the tautological
tuple there. The component test says exactly that this representation
extends to the universal localisation inverting $H$ (use the opposite
convention for the right test). Thus finite-order point tests become
function-field tests for the same matrix. This gives a geometric
criterion, but does not give a terminating equality algorithm for two
upper labels in general. For $A_{n,r,Q}$ it applies to \emph{every}
$Q\in{\rm GL}_r(k)$ in characteristic zero.

\begin{cor}[An explicit equality test for scalar detectors]\label{scalar}
In characteristic zero suppose all chart composition factors of each
of two tame detectors are one-dimensional over $k$. For a detector
with factor set $F$, take all irreducible components of the orbit
closures in $\mathbb A^g_{\bar k}$ of its scalar points, discard
repetitions and keep the inclusion-minimal components. Denote this
finite antichain by $\mathcal C(F)$. Its profile is contained in that
of $F'$ iff each member of $\mathcal C(F')$ contains a member of
$\mathcal C(F)$. In particular the profiles are equal iff these finite
antichains are equal.
\end{cor}
\begin{proof}
On scalar points the determinant of every square matrix is an ordinary
polynomial. Every ordinary polynomial occurs by lifting it, in a fixed
order of monomials, to a one-by-one free polynomial. Passing on a
component means the polynomial does not vanish identically there.
This implication passes from a smaller component to any containing
one. Conversely suppose a component $C'$ on the second list contains
none from the first list. All these components are defined over $k$:
the automorphism permutes the components of an orbit closure; an
orbit point cannot belong to the union of pairwise intersections,
which is closed, invariant and proper. The permutation is transitive
by density. Each component is therefore the closure of one
arithmetic-progression subsequence of $k$-points, hence defined over
$k$ (linear equations in the coefficients of any bounded-degree
vanishing polynomial descend to $k$). For each of the first components
its vanishing ideal consequently cannot contain $I_k(C')$. A finite
union of proper linear subspaces cannot cover the vector space
$I_k(C')$ over the infinite field $k$. Choose $f$ in that ideal
and outside all the first ideals. Its cleared cone passes the first
detector and fails the second. This also treats an empty first list,
for which the zero polynomial can serve. Equality of the two mutual
containment relations on finite antichains is equality of antichains.
\end{proof}

\begin{cor}[A complete invariant in each simple dimension]\label{pure}
Let $k$ be algebraically closed of characteristic zero. Fix $d\ge1$.
Suppose all chart composition factors of both tame detectors have
dimension $d$. Work inside the open set $X_d^{\rm simp}$ of tuples
generating $\operatorname{Mat}_d(k)$. For each orbit component $C$
take the relative closure of
$\mathrm{PGL}_d(C\cap X_d^{\rm simp})$ in this open set, meaning the
closure of the union of simultaneous conjugates. Keep the distinct
inclusion-minimal such varieties. The containment and equality tests of
Corollary~\ref{scalar} hold for these finite antichains of varieties.
Whenever containment fails, a single free \emph{polynomial} supplies a
witness cone.
\end{cor}
\begin{proof}
Every component meets the simple locus densely, since the orbit
of a simple consists of simples and each component receives a dense
subsequence. Its conjugation-saturated relative closure is irreducible
because $\mathrm{PGL}_d$ is irreducible. Invertibility of $H$ is invariant
under conjugacy and open. Passing generically on a component is
therefore equivalent to passing generically on its saturated relative
closure. This proves sufficiency of the containment criterion.

Suppose a target variety $V'$ contains none of the finitely many
source varieties. From each source choose a point $w_i$ outside $V'$.
Let $I\subset R$ be the two-sided ideal of free polynomials that
vanish as whole matrices on $V'$. By
\href{https://arxiv.org/abs/math/0407152}{Reichstein--Vonessen,
Lemma 3.6(b)}, a relatively closed $\mathrm{PGL}_d$-invariant subset
of the generating tuples is the zero set there of its generic-matrix
ideal. The free algebra maps onto the generic matrix ring, so
$I(w_i)\ne0$ for every $i$. Simple evaluation maps onto
$\operatorname{Mat}_d(k)$, and the image of $I$ is a two-sided ideal,
hence is the whole matrix algebra. Discard duplicate conjugacy
classes among the chosen finite set. Their simple annihilators are
distinct maximal two-sided ideals, so $R$ surjects onto the product
of the matrix algebras. The image of $I$ is an ideal surjecting onto
each factor, and therefore contains the identity of the product.
Some $f\in I$ evaluates to the identity at all the chosen points.
It passes generically on every source variety and vanishes identically
on $V'$. Its cleared cone separates the profiles. For an empty source
list the zero polynomial suffices. This proves both assertions.
\end{proof}

Using separate dimensions as disjoint spaces fails for mixed comparisons.
For the infinite-order scalar twist and $g\ge2$, the zero scalar detector's
profile is contained in every tame profile, including those from higher
dimensional simples, as proved in \othernote{parameter_jump}{Failure of
constructibility}. Such inclusions remain covered by Theorem~\ref{geo}.

\section{The one-variable boundary}
\begin{cor}
For a rank-one chart in characteristic zero, finite order gives the
finite-orbit version of the projective-line order. An infinite-order
dilation about a $k$-rational point gives the six-element order of
\othernote{parameter_jump}{Failure of constructibility}. A nonzero
translation gives a four-element chain. These exhaust the possibilities.
\end{cor}
\begin{proof}
An automorphism of $k[U]$ has form $U\mapsto aU+b$. If $a\ne1$
translate the unique fixed point to zero and apply the primary-module
argument of the six-element theorem, which uses only the closed test,
the free resolutions on the chart and the lifts. For $a=1,b\ne0$
every closed point has infinite orbit: a finite set of roots stable
under a nonzero multiple of $b$ would contain an infinite additive
progression in characteristic zero. Hence every nonzero finite module
passes a square polynomial matrix in the tail exactly when its
determinant is a nonzero polynomial. The upper row has top and a single
proper label; the lower row has bottom and the label of all simples.
Finite supports on both sides can meet in only finitely many twists,
so the latter label lies below the proper upper label. The four are
distinct by the zero and nonzero lifts, the source $A$, and $G$.
\end{proof}

For instance the Jordan plane $xz=zx+x^2$, $|x|=|z|=1$, realizes the
translation chart with $h=x$, $U=zx^{-1}$ and $\sigma U=U+1$.
The words $z^i x^j$ form a basis: $xz\mapsto zx+xx$ decreases first
the number of $z$'s or, when equal, the number of inversions. There
is no overlapping leading word. Right multiplication by $x$ is
injective and $xz=(z+x)x$ gives normality by the invertible substitution
$z\mapsto z+x$. The fibre is $k[z]$. A finite right chart module $W$
has a perfect lift with relations $wz-(wU)x$ on a basis of $W$:
highest $z$-degree proves independence, and the Hilbert series shows
that the cokernel maps isomorphically to the nonnegative induced tail.
The opposite ring has the inverse translation and the same proof.
The finite ordinary local dual has reduced bar dual $dx=dz=0$,
$dy=xz-zx-x^2$ with $|y|=1$. Its only radical products land in the
one-dimensional socle. Filtering by the expanded word $y\mapsto xz$
using the same triangular order gives one contractible choice per
disjoint occurrence of $xz$, proving the quotient quasi-isomorphism.

\section{Positive characteristic}
\begin{prop}\label{positive}
For a chart with a linear twist of the generators in characteristic
$p>0$, the set of failing twists for a square $H$ on a given finite
tuple is $p$-automatic. Over an explicitly presented field finitely
generated over $\mathbb F_p$ its finiteness is decidable from the
matrix, tuple and twist. The componentwise geometric test can fail.
\end{prop}
\begin{proof}
Write $J\in{\rm GL}_g(k)$ for the linear substitution. The block
determinant after $j$ substitutions is a polynomial in the entries
of $J^j$. Every monomial of degree $b$ is a matrix coefficient
of $(J^{\otimes b})^j$. A finite direct sum, allowing repeated summands
and the degree-zero constant, expresses the determinant as a linear
functional of powers of one finite matrix. Its characteristic polynomial
therefore gives a scalar linear recurrence; the same holds with
$J^{-1}$. \href{https://arxiv.org/abs/math/0510583}{Derksen's
positive-characteristic Skolem--Mahler--Lech theorem and Theorem 1.10}
give automaticity and the stated effective description. A finite
automaton recognizes infinitely many canonical base-$p$ expansions
exactly when there is a productive cycle of unbounded input length.
This decides the chart test with the data specified above.

Over $k=\mathbb F_p(t)$ take $A_{2,2,Q}$ with
$Q=\operatorname{diag}(t,t+1)$. At the four-generator chart take
the one-dimensional tuple with all entries one. Test
$H=U_{2,0}-U_{1,0}-1$, represented by the cone of $z_2-z_1-x$.
At twists $2j$ and $2j+1$ its value is
\[
                      (t+1)^j-t^j-1 .
\]
It vanishes for all $j=p^a$. In fact these are exactly its zeros:
for $j=p^a m$ with $p\nmid m$, $m>1$, the coefficient of
$t^{p^a}$ is $m$; at $j=0$ the value is $-1$.
But the even iterates are dense in the torus-closure parametrized
by $(u,u,v,v)$, grouping the two slots for each generator. Indeed
the Laurent monomials $t^i(t+1)^j$ are pairwise distinct by the
valuations at $0,-1$. A linear combination vanishing at every power
would give a Vandermonde system with distinct nonzero bases.
The odd closure is the corresponding translate, with the first slot of
each pair still $u,v$ and the other slots $tu,(t+1)v$.
These affine closures are irreducible, and the determinant on both
is $v-u-1$, a nonzero polynomial. These are all components (they
are distinct planes). The failure is consequently a genuine failure
of generic ranks to detect bounded Ext in a two-term dg kernel.
Even the direction of iteration matters: at even negative twists the
scalar is $(t+1)^{-j}-t^{-j}-1$ for $j>0$, which has a pole at
$t=0$ and cannot vanish. At negative odd twists the same statement
holds with the corresponding integer from the two equal slots.
\end{proof}

\paragraph{Priority and boundary.}
The orbit arithmetic is previously established mathematics.
\href{https://arxiv.org/abs/1611.05885}{Bell--Satriano--Sierra}
study fixed homological degree of sheaf Tor as the input moves under
iteration. \href{https://arxiv.org/abs/1706.06470}{Piontkovski}
encodes linear orbits and hyperplane intersections in Hilbert functions
of graded algebras. Here iterating changes the \emph{cohomological
degree of a fixed dg pair}; the result orders all detectors on opposite
module sides. A positive-characteristic test for arbitrary nonlinear
automorphisms beyond the return-time condition is open here.
\end{document}
